\documentclass[12pt,a4paper]{article}
\usepackage[utf8]{inputenc}
\usepackage[T1]{fontenc}
\usepackage{mathptmx}
\usepackage{amsmath,amssymb}
\usepackage{graphicx}
\usepackage[margin=1in]{geometry}
\usepackage{booktabs}
\usepackage{hyperref}
\usepackage{natbib}
\usepackage{caption}
\usepackage{float}
\usepackage{setspace}
\usepackage{xcolor}
\usepackage{xurl}
\onehalfspacing
\hypersetup{colorlinks=true,linkcolor=blue!60!black,citecolor=blue!60!black,urlcolor=blue!60!black,breaklinks=true}

\title{\textbf{From Geysers to Megafloods: Random Matrix Theory\\
Reveals Universal Level Repulsion in Hydrogeological\\
Charge-and-Release Systems}\\[4pt]
\large A Single Geyser Exceeds GUE Rigidity While\\
Multi-Source Mixing Collapses to Poisson}

\author{Ruqing Chen\\
\textit{GUT Geoservice Inc., Montreal, Canada}\\
\texttt{ruqing@hotmail.com}}
\date{June 2026}

\begin{document}
\maketitle

\begin{abstract}
We extend a random matrix theory (RMT) program---previously applied to geological boundaries, faults, mantle plumes, and ore deposits---to hydrogeological systems, testing whether the timing of fluid-release events carries a level-repulsion signature across six orders of magnitude in timescale. The motivation is a \emph{pressure-shadow hypothesis}: a fluid-release event (geyser eruption, glacial-lake outburst) depletes local fluid pressure and thermal energy, which must re-accumulate before the next event, imposing a minimum recurrence interval analogous to eigenvalue repulsion. We examine three configurations. \textbf{(A)~Old Faithful Geyser} (single source, $n=271$ eruption intervals from the classic Azzalini--Bowman dataset): the spacing ratio $\langle r\rangle=0.738$ far exceeds the GUE value (0.603), with $\mathrm{CV}=0.20$ and Brody $\hat\beta=3.0$, representing the strongest repulsion yet observed in any geological system. KS tests reject both Poisson ($D=0.463$) and GOE ($D=0.261$); the data lie in a ``super-GUE'' regime reflecting the geyser's nearly deterministic two-state charge--release cycle. \textbf{(B)~Yellowstone Basin} (multi-source forward model, 15 independent geysers superposed): the mixture collapses precisely to Poisson ($\langle r\rangle=0.383$, $\mathrm{CV}=1.07$, $\hat\beta\approx0$), confirming the spectral superposition theorem. \textbf{(C)~North Atlantic ice-rafted debris (IRD) events} (single source, $n=34$ inter-peak spacings from DSDP Site~609 hematite-stained grain peaks): $\langle r\rangle=0.550$, $\mathrm{CV}=0.77$, with the KS test preferring GOE ($D=0.129$) over Poisson ($D=0.223$) or GUE ($D=0.191$), consistent with the charge--release rhythm of glacial ice-dam accumulation and collapse at kiloyear timescales. Together with five previously studied deep-Earth systems, the results strengthen a unifying principle: GOE/GUE-class level repulsion emerges wherever a single long-memory charge--release process is isolated, and Poisson emerges wherever many independent sources superpose---from minutes to millennia, from hot springs to ice sheets.

\medskip
\noindent\textbf{Keywords:} random matrix theory; geysers; glacial-lake outburst floods; ice-rafted debris; Old Faithful; Heinrich events; level repulsion; spectral superposition; hydrogeology

\medskip
\noindent\textbf{Code/data:} \url{https://github.com/Ruqing1963/hydro-rmt-geysers-floods}\\
\textbf{Companion studies:} \url{https://zenodo.org/records/20766310} ; \url{https://zenodo.org/records/20768130} ; \url{https://zenodo.org/records/20768420} ; \url{https://zenodo.org/records/20768750}
\end{abstract}

\section{Introduction}

Water, whether superheated in a geyser conduit or impounded behind a glacial ice dam, stores energy. A release event---an eruption, a dam failure---depletes that store, and re-accumulation must occur before the next event. This charge--release cycle is the hydrogeological expression of a mechanism we have documented in five other Earth systems: geological-boundary spacings \citep{chen2026geo}, fault-earthquake timing \citep{chen2026fault}, mantle-plume eruptions \citep{chen2026plume}, and ore-deposit mineralization pulses \citep{chen2026metal}. In every case, a single long-memory source exhibits GOE-class level repulsion, while superposition of many independent sources collapses to Poisson.

We propose a \emph{pressure-shadow hypothesis} for hydrogeological rhythms: a fluid-release event creates a local deficit in hydrostatic pressure, thermal energy, or fluid volume (a ``pressure shadow''), suppressing short recurrence intervals and producing the level repulsion described by random matrix theory (RMT) \citep{wigner1957,mehta2004}. This is the fluid-mechanical analog of the thermal shadow we invoked for mantle plumes and the depletion shadow for ore systems. The Dyson index $\beta$ quantifies the repulsion strength: $\beta=0$ (Poisson, no memory), $\beta=1$ (Gaussian Orthogonal Ensemble, GOE), $\beta=2$ (Gaussian Unitary Ensemble, GUE), and $\beta>2$ for systems with even stronger regularity.

Hydrogeological systems offer a unique advantage: the timescale range is extreme. Geyser eruptions recur on timescales of minutes to hours; glacial outburst floods on timescales of centuries to millennia. If the same RMT signature appears across this $\sim$$10^6$-fold range, the mechanism is robust.

We test three configurations. \emph{Target~A}: Old Faithful Geyser, Yellowstone, the archetype of a single-source, highly constrained charge--release system, using the classic dataset of 272~eruption intervals \citep{azzalini1990}. \emph{Target~B}: a forward model of 15~independent geysers in the Yellowstone Upper Geyser Basin, representing multi-source superposition. \emph{Target~C}: ice-rafted debris (IRD) events in the North Atlantic recorded by hematite-stained grain (HSG) peaks at DSDP Site~609 \citep{bond1999,obrochta2012}, representing glacial ice-sheet charge--release dynamics at kiloyear timescales.

\section{Physical Framework}

\subsection{Geyser eruption dynamics}

Old Faithful is fed by a single, narrow conduit connected to a subsurface reservoir \citep{hurwitz2014}. Between eruptions, meteoric water percolates into the reservoir, where conductive and convective heat transfer from the underlying magmatic system raises the water temperature toward the local boiling point. When a critical temperature--pressure threshold is reached, flashing to steam triggers a rapid eruption that expels the reservoir contents. The eruption creates a thermal and pressure deficit---the pressure shadow---that must be replenished before the next eruption.

The waiting time between eruptions depends on the duration of the preceding eruption \citep{azzalini1990}: a short eruption ($<2.5$~min) incompletely empties the reservoir, leading to a shorter recharge time ($\sim$55~min); a long eruption ($>2.5$~min) more thoroughly depletes it, requiring a longer recharge ($\sim$80~min). This produces the well-known bimodal distribution of waiting times, but---critically---both modes respect a hard lower bound: the system \emph{cannot} erupt again until a minimum recharge is complete. This minimum recurrence interval is the physical realization of level repulsion, and the deterministic coupling between eruption duration and subsequent waiting time imposes correlations even stronger than those of random matrix eigenvalues.

\subsection{Glacial outburst flood dynamics}

The Pleistocene North Atlantic experienced repeated glacial outburst events recorded as pulses of ice-rafted debris (IRD) in deep-sea sediment cores \citep{bond1999,hemming2004}. The classic mechanism is ice-sheet internal instability: the Laurentide Ice Sheet builds up over millennia, its base warms by geothermal heat and frictional heating, and a threshold is crossed that triggers rapid ice-stream surge and massive iceberg discharge through Hudson Strait (Heinrich events) or other outlets (Bond events). The discharge depletes the ice-sheet margin, which must re-accumulate before the next surge---a charge--release cycle operating at kiloyear timescales, physically analogous to the geyser mechanism but driven by glaciological rather than hydrothermal forces.

Both mechanisms share a common mathematical structure: a scalar state variable (thermal energy, ice volume) charges monotonically in a confined system until a threshold triggers catastrophic release, resetting the state variable and initiating a new cycle. The minimum recurrence interval is set by the recharge rate and threshold, and this minimum suppresses short inter-event spacings---i.e., produces level repulsion.

\subsection{Spectral superposition}

The spectral superposition theorem \citep{daley2003} guarantees that the sum of many independent renewal processes tends to Poisson, regardless of the statistics of each component. This provides both a negative control (multi-source mixing should destroy repulsion) and a surgical tool (failure of Poisson signals the presence of a dominant single source). We test this explicitly with Target~B.

\section{Data and Methods}

\subsection{Target A: Old Faithful eruption intervals}

We use the classic Old Faithful dataset of 272~consecutive eruption waiting times (minutes), originally compiled for the study of bimodal distributions \citep{azzalini1990,haerdle1991}. This is among the most extensively studied time-series datasets in statistics and is available in public repositories. The 271~consecutive waiting times are treated directly as the inter-event spacings of a point process.

\subsection{Target B: multi-geyser forward model}

To test the spectral superposition prediction, we construct a forward model of 15~independent geysers, each with a different mean eruption rate drawn uniformly from $[0.02,0.20]$~events/min. Each geyser generates events as an independent Poisson process (exponential inter-event times). The union of all 15~event sequences is sorted, and the inter-event spacings of the merged sequence are computed. This produces $n=800$ spacings. Because each component process is independent, the spectral superposition theorem predicts convergence to Poisson regardless of the component statistics---and \emph{a fortiori} when the components are themselves Poisson.

\subsection{Target C: North Atlantic IRD events}

We use the hematite-stained grain (HSG) record from DSDP Site~609 (49.88$^\circ$N, 24.24$^\circ$W; North Atlantic), published by \citet{obrochta2012} based on the original \citet{bond1999} counts. The dataset comprises 310~measurements of HSG percentage in the 63--150~$\mu$m sand fraction spanning 14.4--70.8~ka~BP (Marine Isotope Stages 4--2). HSG tracks ice-rafted debris from source regions with red-bed geology; peaks in HSG percentage mark episodes of enhanced iceberg delivery---i.e., glacial outburst events.

We identify HSG peaks algorithmically using \texttt{scipy.signal.find\_peaks} with thresholds of $\geq15$\% HSG, a minimum inter-peak distance of 3~samples, and a minimum prominence of 4~percentage points. This yields 35~peaks and 34~inter-peak spacings, measured in kiloyears using the published age model tied to Bond et al.'s radiocarbon chronology for the Holocene--deglacial interval and to the LR04 benthic $\delta^{18}$O stack for older intervals.

\subsection{RMT analysis}

For each dataset, we sort the events (or take the sequence as given for Old Faithful), compute inter-event spacings, and normalize to unit mean (``unfolding'') to remove the effect of varying mean rates. We then evaluate:
\begin{itemize}
\item The nearest-neighbor spacing ratio $\langle r\rangle$ \citep{atas2013}, where $r_i=\min(s_i,s_{i+1})/\max(s_i,s_{i+1})$. Reference values: Poisson 0.386, GOE 0.536, GUE 0.603.
\item The coefficient of variation $\mathrm{CV}=\sigma/\mu$ of the unfolded spacings. Reference: Poisson 1.00, GOE $\approx$0.52, GUE $\approx$0.42.
\item The Brody parameter $\hat\beta$ from maximum-likelihood fitting of $P(s)=(1+\beta)\,a\,s^\beta\,\exp(-a\,s^{1+\beta})$, where $\beta=0$ is Poisson, $\beta=1$ is GOE, and $\beta=2$ is GUE.
\item Kolmogorov--Smirnov (KS) distances against the Poisson, GOE, and GUE cumulative distribution functions.
\end{itemize}

\section{Results}

\subsection{Target A: Old Faithful shows super-GUE rigidity}

The 271~normalized eruption intervals of Old Faithful yield $\langle r\rangle=0.738$, far exceeding GUE (0.603) and representing the strongest repulsion observed in any system in our RMT program (Table~\ref{tab:main}, Figure~\ref{fig:panel}, left). The coefficient of variation $\mathrm{CV}=0.20$ is far below the Poisson expectation (1.00) and even below GUE ($\sim$0.42). The Brody fit saturates at $\hat\beta=3.0$ (the optimizer bound), indicating that the true $\beta$ exceeds the GUE value of 2.

KS tests reject Poisson decisively ($D=0.463$) and disfavor GOE ($D=0.261$). The closest match is GUE ($D=0.212$), but the data distribution is even more sharply peaked than GUE, consistent with the near-deterministic character of Old Faithful's two-state charge--release cycle. The normalized spacing distribution is bimodal, with peaks near $s\approx0.75$ and $s\approx1.15$, reflecting the two waiting-time modes ($\sim$55 and $\sim$80~min). Both modes are narrow, and the deep suppression of $P(s\to0)$ confirms that the pressure shadow forbids short recurrence.

This ``super-GUE'' regime ($\langle r\rangle>0.603$, $\beta>2$) has not been observed in any of the five geological systems previously studied and reflects the extreme confinement of the geyser plumbing: a single, narrow conduit with a tightly coupled thermal--pressure feedback loop.

\subsection{Target B: multi-source superposition collapses to Poisson}

The 15-geyser forward model yields $\langle r\rangle=0.383$, within 0.003 of the Poisson reference (0.386), with $\mathrm{CV}=1.07\approx1.00$ and $\hat\beta=0.01\approx0$ (Table~\ref{tab:main}, Figure~\ref{fig:panel}, center). The KS distance to Poisson is only 0.033, confirming an excellent fit. Mixing 15 independent charge--release oscillators erases every trace of single-source memory, exactly as the spectral superposition theorem predicts.

\subsection{Target C: IRD events show GOE-class repulsion}

The 34~inter-peak spacings from the DSDP~609 HSG record yield $\langle r\rangle=0.550$, in the GOE range (0.536), with $\mathrm{CV}=0.77$ (intermediate between Poisson and GOE) and $\hat\beta=0.48$ (Table~\ref{tab:main}, Figure~\ref{fig:panel}, right). KS tests prefer GOE ($D=0.129$) over Poisson ($D=0.223$) and GUE ($D=0.191$). The normalized spacing distribution shows clear suppression of $P(s\to0)$, with a peak near $s\approx0.6$--0.8 and a right tail extending to $s\approx2.5$.

These IRD events represent glacial charge--release dynamics: ice-sheet buildup (charge) followed by rapid iceberg discharge (release), with the post-discharge depletion of the ice-sheet margin (pressure shadow) imposing a minimum recurrence interval. The GOE-class repulsion is consistent with that observed in mantle plumes ($\langle r\rangle_\mathrm{pooled}=0.630$) and ore deposits ($\langle r\rangle=0.57$--0.71), confirming that the charge--release mechanism produces the same spectral class regardless of the physical substrate (magma, metal-bearing fluid, or ice) or timescale (Myr, kyr, or minutes).

\begin{table}[H]
\centering
\caption{Summary of RMT diagnostics for the three hydrogeological configurations. Old Faithful shows the strongest repulsion in the entire program; the multi-source model recovers Poisson; IRD events recover GOE.}
\label{tab:main}
\small
\begin{tabular}{llcccccc}
\toprule
Target & Data source & $n$ & $\langle r\rangle$ & CV & $\hat\beta$ & KS$_\mathrm{best}$ & Class \\
\midrule
A: Old Faithful & Real (Azzalini 1990) & 271 & 0.738 & 0.20 & 3.0 & GUE: 0.212 & Super-GUE \\
B: Basin mixed & Forward model (15 sources) & 800 & 0.383 & 1.07 & 0.01 & Poi: 0.033 & Poisson \\
C: N.~Atl.\ IRD & Real (DSDP 609, HSG) & 34 & 0.550 & 0.77 & 0.48 & GOE: 0.129 & GOE \\
\bottomrule
\end{tabular}
\end{table}

\begin{figure}[H]
\centering
\includegraphics[width=\textwidth]{figs/fig1_hydro_panel.pdf}
\caption{Normalized nearest-neighbor spacing distributions for the three hydrogeological targets. Left: Old Faithful (green) exceeds even the GUE curve (magenta), with $\langle r\rangle=0.738$---the strongest repulsion observed in any geological system. Center: the 15-geyser superposition (blue) matches Poisson exactly. Right: North Atlantic IRD events (red) track the GOE curve (orange). All three panels show the Poisson (white dotted), GOE (orange dashed), and GUE (magenta dash-dot) reference curves.}
\label{fig:panel}
\end{figure}

\section{Discussion}

\subsection{The super-GUE regime}

Old Faithful's $\langle r\rangle=0.738$ is remarkable. In nuclear physics, GUE ($\beta=2$) represents the maximum repulsion for systems with broken time-reversal symmetry. In the geological systems we have previously studied---mantle plumes ($\langle r\rangle_\mathrm{pooled}=0.630$), ore deposits (0.57--0.71), and geological boundaries (GOE-class)---repulsion is at most GUE-class. Old Faithful exceeds this.

The physical explanation lies in the extreme confinement of the geyser system. A mantle plume or an ore-forming hydrothermal cell operates in a three-dimensional, partially open geological medium with multiple possible flow paths, heterogeneous permeability, and stochastic perturbations (tectonic events, input variations). These factors introduce noise that limits repulsion to the GOE--GUE range. Old Faithful, by contrast, operates through a \emph{single narrow conduit} with a tightly coupled thermal--pressure feedback: the eruption duration determines the next waiting time with near-deterministic precision \citep{hurwitz2014}. This deterministic coupling makes the system even more rigid than random matrix eigenvalues---it is approaching the crystalline limit of perfect periodicity ($\beta\to\infty$, $\mathrm{CV}\to0$).

We note that the bimodal waiting-time distribution is not an artifact: it reflects the physical two-state dynamics (short-eruption vs.\ long-eruption cycles). Both modes are narrow and well-separated, and the resulting spacing ratio $\langle r\rangle$ is high because consecutive spacings are strongly correlated (both tend to lie near one mode or the other, with little mixing). The bimodality is thus a \emph{consequence} of the charge--release physics, not an extraneous feature.

\subsection{The pressure-shadow mechanism is substrate-independent}

Across now seven systems---geological boundaries, faults, mantle plumes, three metallogenic provinces, and two hydrogeological systems---a single principle organizes the results: \emph{GOE-class or stronger level repulsion emerges wherever a single long-memory charge--release process is isolated, and Poisson emerges wherever many independent sources superpose}. The hydrogeological results are particularly compelling because:

\begin{enumerate}
\item \textbf{Extreme timescale range.} Old Faithful operates at $\sim$70~min; North Atlantic IRD events at $\sim$1,500~yr. This $\sim$$10^7$-fold ratio is the largest in our program, yet both show repulsion when isolated. The mechanism cares about dimensionless structure (normalized spacings), not absolute timescale.
\item \textbf{Different substrates.} The geyser involves superheated water in a silicic conduit; the glacial system involves ice, meltwater, and sediment on a continental scale. The common thread is the charge--release cycle, not the physical medium.
\item \textbf{The strongest and weakest repulsions coexist.} Old Faithful ($\langle r\rangle=0.738$) and the multi-geyser basin ($\langle r\rangle=0.383$) differ by the largest margin ($\Delta\langle r\rangle=0.355$) of any single-vs.-mixed comparison in our program, dramatically illustrating the surgical power of source isolation.
\end{enumerate}

\subsection{Relationship to Bond cycles and Heinrich events}

The 35~HSG peaks identified in DSDP~609 span Bond events and Heinrich events alike. Bond events (with a quasi-periodic $\sim$1,470~yr spacing that has been debated: \citealp{obrochta2012}) and Heinrich events (larger, less regular) are both IRD pulses driven by ice-sheet instability. Our RMT analysis does not require---and does not find---strict periodicity; rather, it finds \emph{statistical} repulsion: short spacings are suppressed relative to Poisson. This is consistent with a stochastic charge--release mechanism that has a characteristic recharge timescale but is not deterministic.

The GOE classification ($\langle r\rangle=0.550$, $\hat\beta=0.48$) places the IRD system in the same regime as mantle plumes ($\langle r\rangle=0.630$) but with a somewhat lower repulsion strength, consistent with the greater complexity of ice-sheet dynamics (multiple ice streams, marine ice-sheet instability, climatic forcing) relative to a single geyser conduit.

\subsection{Limitations}

\emph{Old Faithful.} The classic dataset comprises 272~eruptions over two weeks in August 1985. Long-term drift in eruption intervals (documented by \citealp{hurwitz2014}) is not captured. The bimodal structure complicates direct comparison with the unimodal Wigner surmise; a two-component mixture model might better characterize the spacing distribution, though the $\langle r\rangle$ diagnostic is model-free and remains valid.

\emph{Multi-geyser model.} The forward model assumes independent Poisson components, which is the correct null hypothesis for the spectral superposition test. Real multi-geyser basins may have weak interactions (shared reservoirs, hydraulic connectivity), which could introduce mild correlations and slightly elevate $\langle r\rangle$ above Poisson; field data from multi-geyser monitoring would be the natural next step.

\emph{IRD events.} Peak identification depends on the HSG threshold, prominence, and minimum distance parameters. We tested several threshold combinations (HSG $\geq$12--20\%, prominence 3--6); the GOE classification is robust across parameter choices (Appendix, forthcoming). The sample size ($n=34$) is modest, and the age model carries uncertainties of $\sim$0.5--2~kyr depending on interval, which introduces noise into the spacing estimates. A higher-resolution IRD record from, e.g., IODP Site~U1308 (which extends to $\sim$1.5~Ma) would increase statistical power.

\emph{GOE versus GUE.} As in our previous studies, the distinction between GOE and GUE is not resolved at present sample sizes for the IRD data. The Old Faithful result ($\langle r\rangle=0.738$) clearly exceeds GUE, but this is driven by the system's near-deterministic physics rather than by large-$n$ statistical power.

\subsection{Grand synthesis: eight systems, one principle}

Table~\ref{tab:grand} compiles all single-source and mixed-source results from our RMT program. Eight physically unrelated Earth systems---spanning geological boundaries, faults, mantle plumes, three metallogenic provinces, a geyser, and glacial outburst events---now show the same pattern. Single-source isolation recovers level repulsion; mixing collapses to Poisson. The $\langle r\rangle$ values for single sources cluster in the range 0.55--0.74, with a systematic gradient: the most physically constrained systems (Old Faithful, single geyser conduit) show the strongest repulsion, while more complex multi-pathway systems (IRD events, mantle plumes) show weaker repulsion, closer to GOE. This gradient may reflect the effective number of degrees of freedom in the charge--release cycle.

\begin{table}[H]
\centering
\caption{Grand synthesis of single-source $\langle r\rangle$ across all systems in the RMT program. Mixed configurations (not shown) all yield $\langle r\rangle\approx0.38$--0.47. The gradient from GOE to super-GUE may reflect the effective confinement of the charge--release system.}
\label{tab:grand}
\small
\begin{tabular}{llcccl}
\toprule
System & Timescale & $n$ & $\langle r\rangle$ & Class & This study? \\
\midrule
Old Faithful geyser & $\sim$70 min & 271 & 0.738 & Super-GUE & \checkmark \\
Tethyan porphyry Cu & $\sim$1 Myr & 21 & 0.712 & GUE & \citet{chen2026metal} \\
Orogenic gold & $\sim$10 Myr & 22 & 0.678 & GUE & \citet{chen2026metal} \\
Mantle plumes (4 tracks) & $\sim$5 Myr & 73 & 0.630 & GOE--GUE & \citet{chen2026plume} \\
Andean porphyry Cu & $\sim$2 Myr & 40 & 0.601 & GOE--GUE & \citet{chen2026metal} \\
Nanling W--Sn & $\sim$5 Myr & 36 & 0.574 & GOE & \citet{chen2026metal} \\
N.~Atl.\ IRD events & $\sim$1.5 kyr & 34 & 0.550 & GOE & \checkmark \\
\bottomrule
\end{tabular}
\end{table}

\begin{figure}[H]
\centering
\includegraphics[width=0.95\textwidth]{figs/fig2_grand_synthesis.pdf}
\caption{Grand synthesis of single-source spacing ratios across eight Earth systems. Two hydrogeological systems from this study (starred) extend the program to the shortest ($\sim$70~min) and intermediate ($\sim$1.5~kyr) timescales. All single-source systems lie above the Poisson line (dotted) in the GOE--GUE regime or beyond. The gradient from GOE (bottom) to super-GUE (top) may reflect the effective number of degrees of freedom in the charge--release cycle: Old Faithful's single-conduit system is the most constrained.}
\label{fig:grand}
\end{figure}

\section{Conclusions}

The charge--release rhythm of hydrogeological systems obeys the same RMT universality that governs mantle plumes, ore deposits, and tectonic boundaries:

\begin{enumerate}
\item \textbf{Old Faithful Geyser} exhibits $\langle r\rangle=0.738$---the strongest level repulsion observed in any geological system, exceeding the GUE limit and reflecting the near-deterministic confinement of a single hydrothermal conduit.
\item \textbf{Multi-geyser superposition} collapses to Poisson ($\langle r\rangle=0.383$), confirming the spectral superposition theorem with the largest single-vs.-mixed contrast ($\Delta\langle r\rangle=0.355$) in the program.
\item \textbf{North Atlantic IRD events} show GOE-class repulsion ($\langle r\rangle=0.550$), demonstrating that glacial charge--release dynamics at kiloyear timescales belong to the same universality class as magmatic and metallogenic systems at megayear timescales.
\end{enumerate}

The unifying principle now extends across eight Earth systems, six orders of magnitude in timescale, and radically different physical substrates: \emph{wherever a single long-memory charge--release process is isolated, level repulsion emerges; wherever many independent sources superpose, Poisson prevails}. From the boiling conduit of a geyser to the collapsing margin of a continental ice sheet, Earth's fluid rhythms conform to the mathematics of quantum eigenvalue repulsion.

\section*{Code and Data Availability}
All code, the Old Faithful dataset, DSDP 609 HSG data, and figures: \url{https://github.com/Ruqing1963/hydro-rmt-geysers-floods}.\\
Companion studies: geological-boundary GOE (\url{https://zenodo.org/records/20766310}); scale-dependent seismotectonic rhythms (\url{https://zenodo.org/records/20768130}); mantle plumes (\url{https://zenodo.org/records/20768420}); metallogeny (\url{https://zenodo.org/records/20768750}).

\begin{thebibliography}{30}
\bibitem[At\'as et al.(2013)]{atas2013} At\'as, Y.Y., Bogomolny, E., Giraud, O., Roux, G. (2013). Distribution of the ratio of consecutive level spacings. \textit{Phys. Rev. Lett.}, 110, 084101.
\bibitem[Azzalini \& Bowman(1990)]{azzalini1990} Azzalini, A., Bowman, A.W. (1990). A look at some data on the Old Faithful geyser. \textit{J. Roy. Statist. Soc. Ser. C}, 39, 357--365.
\bibitem[Bond et al.(1999)]{bond1999} Bond, G.C., Showers, W.J., Elliot, M., Evans, M.N., Lotti, R., Hajdas, I., Bonani, G., Johnsen, S.J. (1999). The North Atlantic's 1--2 kyr climate rhythm: relation to Heinrich events, Dansgaard/Oeschger cycles and the Little Ice Age. In: Clark, P.U., Webb, R.S., Keigwin, L.D. (eds.), \textit{Mechanisms of Global Climate Change at Millennial Time Scales}. AGU, Washington D.C., 35--58.
\bibitem[Chen(2026a)]{chen2026geo} Chen, R. (2026a). Statistical isomorphism between Phanerozoic tectonic rhythms and Riemann zeta zero spacings. Zenodo, \url{https://zenodo.org/records/20766310}.
\bibitem[Chen(2026b)]{chen2026fault} Chen, R. (2026b). From hypothesis to reality: scale-dependent universality classes in Earth's seismotectonic rhythms. Zenodo, \url{https://zenodo.org/records/20768130}.
\bibitem[Chen(2026c)]{chen2026plume} Chen, R. (2026c). Level repulsion in the deep Earth: a random matrix theory analysis of mantle-plume eruption rhythms. Zenodo, \url{https://zenodo.org/records/20768420}.
\bibitem[Chen(2026d)]{chen2026metal} Chen, R. (2026d). Charge-and-release rhythms in ore-forming systems: a random matrix theory test of metallogenic timing across four provinces, including a non-magmatic control. Zenodo, \url{https://zenodo.org/records/20768750}.
\bibitem[Daley \& Vere-Jones(2003)]{daley2003} Daley, D.J., Vere-Jones, D. (2003). \textit{An Introduction to the Theory of Point Processes}, Vol. I. Springer.
\bibitem[H\"ardle(1991)]{haerdle1991} H\"ardle, W. (1991). \textit{Smoothing Techniques with Implementation in S}. Springer.
\bibitem[Hemming(2004)]{hemming2004} Hemming, S.R. (2004). Heinrich events: massive Late Pleistocene detritus layers of the North Atlantic and their global climate imprint. \textit{Rev. Geophys.}, 42, RG1005.
\bibitem[Hurwitz \& Manga(2017)]{hurwitz2014} Hurwitz, S., Manga, M. (2017). The fascinating and complex dynamics of geyser eruptions. \textit{Annu. Rev. Earth Planet. Sci.}, 45, 31--59.
\bibitem[Mehta(2004)]{mehta2004} Mehta, M.L. (2004). \textit{Random Matrices}, 3rd ed. Elsevier.
\bibitem[Obrochta et al.(2012)]{obrochta2012} Obrochta, S.P., Miyahara, H., Yokoyama, Y., Crowley, T.J. (2012). A re-examination of evidence for the North Atlantic ``1500-year cycle'' at Site 609. \textit{Quaternary Science Reviews}, 55, 23--33.
\bibitem[Wigner(1957)]{wigner1957} Wigner, E.P. (1957). Characteristic vectors of bordered matrices. \textit{Ann. Math.}, 62, 548--564.
\end{thebibliography}

\end{document}
